% Generated draft: requires mathematical and source review.
\documentclass[11pt]{article}
\usepackage[margin=25mm]{geometry}
\usepackage{fontspec}
\setmainfont{Latin Modern Roman}
\usepackage{amsmath,amssymb,mathtools,unicode-math}
\setmathfont{Latin Modern Math}
\usepackage{enumitem,needspace,xurl,hyperref}
\setlength{\parindent}{0pt}
\setlength{\parskip}{5pt}
\setlength{\emergencystretch}{3em}
\newcommand{\N}{\mathbb{N}}
\newcommand{\Z}{\mathbb{Z}}
\newcommand{\Q}{\mathbb{Q}}
\newcommand{\R}{\mathbb{R}}
\newcommand{\C}{\mathbb{C}}
\newcommand{\F}{\mathbb{F}}
\newcommand{\A}{\mathbb{A}}
\newcommand{\PP}{\mathbb{P}}
\newcommand{\E}{\mathbb{E}}
\newcommand{\eps}{\varepsilon}
\newcommand{\dd}{\,\mathrm{d}}
\begin{document}
\section*{UnsolvedMath ID: 30005190}
\textbf{Source problem code:} catalogue entry
\subsection{Definitions and mathematical statement}
Let \(n\) be a positive integer with \(n-1\neq 4\), and let \(M\) be a closed connected smooth manifold of dimension \(n-1\). Here closed means compact without boundary. Assume that there is no smooth Riemannian metric \(h\) on \(M\) satisfying \(\operatorname{Scal}_h(x)>0\) for every \(x\in M\). Set
\[
V=M\times[-1,1],\qquad \partial_-V=M\times\{-1\},\qquad \partial_+V=M\times\{1\}.
\]
For a smooth Riemannian metric \(g\) on \(V\), let \(d_g\) be its intrinsic Riemannian distance, defined using lengths of curves contained in \(V\), and define
\[
\operatorname{width}(V,g)=\operatorname{dist}_g(\partial_-V,\partial_+V)
=\inf\{d_g(x,y):x\in\partial_-V,\ y\in\partial_+V\}.
\]
The supplied formulation of Gromov's scalar curvature band-width conjecture asks whether, for every such \(M\) and every such \(g\),
\[
\bigl(\operatorname{Scal}_g(p)\geq n(n-1)\text{ for every }p\in V\bigr)
\quad\Longrightarrow\quad
\operatorname{width}(V,g)\leq\frac{2\pi}{n}.
\]
No orientability hypothesis, product-metric condition, or boundary mean-curvature condition is stated in the supplied conjecture.

\textit{Unresolved qualification.} The record gives no lower dimension bound beyond the dimensional notation and the exclusion \(n-1\neq4\). In particular, it does not explicitly exclude \(n=1\). The intended lower dimension range requires clarification; the formulation above retains the literal range rather than adding a hypothesis.
\subsection{Short English statement}
The supplied conjecture asks whether every metric on the band \(M\times[-1,1]\) with scalar curvature at least \(n(n-1)\) has boundary-to-boundary distance at most \(2\pi/n\), provided \(M\) admits no metric of everywhere positive scalar curvature and \(\dim M\neq4\). The intended lower dimension range is unresolved. This is a statement-only draft requiring human mathematical review.
\subsection{Sources}
\begin{itemize}
\item Local UnsolvedMath registry, supplied complete record: permanent ID \(30005190\); source problem code \texttt{catalogue entry}; title: Gromov's Scalar Curvature Band-Width Conjecture; collection: Oberwolfach Reports Open Problems. The record labels the statement \texttt{corrected\_verified} and says it was checked against the official report. These are supplied verification claims, not independent verification here.
\item Supplied source citation: \emph{Non-Commutative Geometry and Cyclic Homology (2023)}, Oberwolfach Reports, DOI 10.4171/owr/2022/35, pp. 1987--2058. Original Oberwolfach report: \url{https://doi.org/10.4171/owr/2022/35}. The registry records proposed year 2022. The supplied excerpt attributes the conjecture to Gromov and identifies the contribution as \emph{Metric inequalities with scalar curvature via the Dirac operator}, Rudolf Zeidler, joint work with Simone Cecchini. It cites the conjecture as reference [4], Section 11.12, Conjecture C, and mentions references [3, 4]; the corresponding bibliography entries are not supplied.
\item Supplied literature source: Rade, \emph{Scalar and mean curvature comparison via mu-bubbles}, \url{https://arxiv.org/abs/2104.10120}.
\item Supplied literature source: Cecchini, \emph{Nonnegative scalar curvature on manifolds with at least two ends}, \url{https://londmathsoc.onlinelibrary.wiley.com/doi/full/10.1112/topo.12303}.
\item Supplied literature assessment dated exactly \texttt{2026-08-22}: status \texttt{partially\_solved}. The assessment states that the inequality is proved for broad classes, in particular for orientable bands in dimensions \(n\leq7\), excluding \(n=5\), and that the all-dimensional formulation in the record is not settled. This is supplied dated curation derived from the OWR collection, not independently verified here; the source explicitly requests independent verification of time-sensitive status and cited claims.
\end{itemize}
\end{document}
